\chapter*{使用说明}
\addcontentsline{toc}{chapter}{使用说明}
本讲义围绕一个问题展开：给定像素、文本单元（token，可以是一个字、子词或符号）、原子坐标、自旋构型或量子态系数等高维对象，如何表达、学习并采样一个合理的概率分布？生成模型在这里指的是一整套方法：把概率建模、统计推断、数值采样和系统实现连接起来，回答同一个对象空间上的分布问题。

\begin{keybox}[本讲义的四个视角]
每个模型都可以从下面四个视角去看，后文讨论任何一个新方法时，都回到这四个问题：
\begin{enumerate}
  \item \textbf{概率对象：}模型表示的是显式密度、条件概率、隐式分布、对数密度梯度还是能量？
  \item \textbf{统计目标：}训练在做最大似然估计、变分下界、对抗匹配、对数密度梯度匹配，还是约束分布变化幅度的奖励优化？
  \item \textbf{计算路径：}采样是一次前向、逐 token 解码、可逆变换、随机微分方程积分，还是马尔可夫链蒙特卡洛？
  \item \textbf{系统瓶颈：}真正限制规模的是浮点运算次数、显存容量、显存带宽、跨卡通信、随机步数，还是数据质量？
\end{enumerate}
\end{keybox}
